%NO MODIFICAR ESTA SECCION!
\documentclass{article} % Define la clase del documento, en este caso, un artículo

\usepackage[letterpaper,margin=3cm]{geometry} % Configura el tamaño del papel y los márgenes del documento
\usepackage{graphicx} % Permite la inserción de imágenes
%\usepackage[spanish]{babel}% Activar esta configuración para informes en español, ajusta el idioma del documento
\usepackage[usenames]{color} % Permite el uso de colores definidos por nombre en el documento
\usepackage{hyperref} % Habilita enlaces y referencias dentro del documento
\hypersetup{colorlinks=true, linkcolor = black, citecolor= black} % Configura el color de los enlaces y citas
\usepackage{booktabs} % Proporciona comandos para crear tablas de alta calidad
\usepackage{natbib} % Permite el uso de citas y referencias bibliográficas con diferentes estilos
\usepackage{tikz} % Permite la creación de gráficos y diagramas vectoriales directamente en LaTeX
\usepackage{float} % Para controlar la posición de los elementos flotantes, como imágenes, con la opción [H]
\bibliographystyle{agsm} % Define el estilo de citas y bibliografía (en este caso, el estilo AGSM)
\usepackage{diagbox} % Permite crear celdas con líneas diagonales en tablas
\usepackage{listings} % Permite la inclusión y formateo de código fuente en el documento
\usepackage{xcolor} % Paquete para definir y usar colores en el documento
\usepackage{parskip} % Añade espacio entre párrafos en lugar de sangrías
\usepackage{fancyhdr} % Permite personalizar encabezados y pies de página
\usepackage{amsmath} % Proporciona una amplia variedad de entornos y comandos matemáticos

\pagestyle{fancy} % Usa el estilo fancyhdr
\fancyhf{} % Borra todos los encabezados y pies de página
\renewcommand{\headrulewidth}{0pt}
\renewcommand{\footrulewidth}{0pt} % Desactiva la línea horizontal predeterminada en el pie
\setlength{\headheight}{2cm} % Ajusta la altura del encabezado para hacer espacio para la línea
\fancyhead[L]{\raisebox{0.20cm}{\textbf{Fluid Mechanics}}} % Añade el texto en la parte izquierda del encabezado, subiéndolo ligeramente
\fancyhead[R]{\raisebox{0.1cm}{\includegraphics[width=0.25\linewidth]{LOGO_UNIVERSIDAD.jpg}}} % Añade la imagen en la parte derecha del encabezado y súbela un poco
\fancyhead[C]{\rule{\textwidth}{0.6pt}} % Añade una línea horizontal superior centrada
\fancyfoot[C]{\rule{\textwidth}{0.6pt}} % Añade una línea horizontal en el pie de página centrada
\fancyfoot[R]{\raisebox{-1.5\baselineskip}{\thepage}} % Coloca el número de página a la derecha, con suficiente espacio debajo de la línea
\geometry{top=3cm, bottom=2.5cm} % Ajusta los márgenes superior e inferior

% Definición de colores al estilo Visual Studio Code
\definecolor{codegreen}{rgb}{0.25,0.49,0.48} % Comentarios
\definecolor{codegray}{rgb}{0.5,0.5,0.5} % Números y anotaciones
\definecolor{codepurple}{rgb}{0.58,0,0.82} % Palabras clave
\definecolor{backcolour}{rgb}{0.95,0.95,0.92} % Color de fondo

% Configuración del estilo de las celdas de código
\lstset{
    backgroundcolor=\color{backcolour},   % color de fondo; necesita que el paquete color o xcolor esté cargado
    commentstyle=\color{codegreen},       % estilo de comentarios
    keywordstyle=\color{codepurple},      % estilo de palabras clave
    numberstyle=\tiny\color{codegray},    % estilo de los números de línea
    stringstyle=\color{red},              % estilo de las cadenas de texto
    basicstyle=\ttfamily\small,           % estilo del texto básico
    breakatwhitespace=false,              % ajustes de líneas sólo en espacios en blanco
    breaklines=true,                      % ajustar las líneas si son muy largas
    captionpos=b,                         % posición de la leyenda (abajo)
    keepspaces=true,                      % preserva los espacios en el texto; útil si se usa monoespaciado
    numbers=left,                         % dónde poner los números de línea
    numbersep=5pt,                        % qué tan lejos están los números de línea del código
    showspaces=false,                     % mostrar espacios con subrayados particulares; reemplaza 'showstringspaces'
    showstringspaces=false,               % subrayar los espacios dentro de las cadenas solo
    showtabs=false,                       % mostrar tabulaciones en el código con subrayados particulares
    tabsize=2,                            % tamaños de tabulación a 2 espacios
    language=TeX,                         % lenguaje del código
    morecomment=[l]\#,                    % reconocer # como inicio de comentario en Python
    frame=single,                         % agregar un marco simple alrededor del código
    rulecolor=\color{black}               % color del marco
}

\begin{document}
%----------------------------------------------------------------------------------------
%   PORTADA
%Modificar desde aqui en adelante
%----------------------------------------------------------------------------------------
\begin{titlepage}%Inicio de la carátula, solo modificar los datos necesarios
\newcommand{\HRule}{\rule{\linewidth}{0.5mm}} 
\center 
%----------------------------------------------------------------------------------------
%	ENCABEZADO
%----------------------------------------------------------------------------------------
\includegraphics[width=10cm]{LOGO_UNIVERSIDAD.jpg}\\ % Si esta plantilla se copio correctamente, va a llevar la imagen del logo de la facultad.OBS: Es necesario incluir el paquete: graphicx
\vspace{3cm}
%----------------------------------------------------------------------------------------
%	SECCION DEL TITULO
%----------------------------------------------------------------------------------------
\HRule \\[0.4cm]
{ \huge \bfseries Tittle}\\[0.4cm] % Titulo del documento
{ \huge \bfseries Fluid Mechanics}\\[0.4cm] % Titulo del documento
\HRule \\[1.5cm]
 \vspace{5cm}
%----------------------------------------------------------------------------------------
%	SECCION DEL AUTOR
%----------------------------------------------------------------------------------------
\begin{flushright}
    { \textbf{Profesors:}\\
    Patricio Moreno\\
    Sebastian Sepulveda\\
    \vspace{0.2cm}
    \textbf{Assistant:}\\
    Lukas Wolff\\
    \vspace{0.2cm}
    \textbf{Author:}\\
    Your name\\
}
\end{flushright}
\vspace{1cm}
%----------------------------------------------------------------------------------------
%	SECCION DE LA FECHA
%----------------------------------------------------------------------------------------
{\large \textbf{\today}}\\[2cm] % El comando \today coloca la fecha del dia, y esto se actualiza con cada compilacion, en caso de querer tener una fecha estatica, reemplazar el \today por la fecha deseada
\end{titlepage}
%----------------------------------------------------------------------------------------
%  INDICE
%----------------------------------------------------------------------------------------
\newpage
\tableofcontents
\thispagestyle{plain} % Deshabilita el encabezado en la página del índice
\thispagestyle{empty} % Deshabilita el número de página en la página del índice
\newpage

%Se puede agregar un indice de figuras si es nesesario
%\newpage
%\listoffigures 
%\thispagestyle{plain} % Deshabilita el encabezado en la página del índice %
%\thispagestyle{empty}
%\newpage
%----------------------------------------------------------------------------------------
%   ACÁ EMPIEZA EL INFORME
\setcounter{page}{1} % Reinicia el contador de páginas
%----------------------------------------------------------------------------------------
%Este es el formato a seguir para los titulos de las secciones
\section{INTRODUCTION}

\section{PROCEDURE}

\section{THEORY}

\section{ANALYSIS}

\section{QUESTIONS}

\section{CONCLUSIONS}

\section{REFERENCES}

%----------------------------------------------------------------------------------------
%Aquí van unos tips básicos de LaTeX
\newpage
\section{LaTeX Tips}

\subsection{Algunos tips básicos :)}
\textbf{Texto en negrita.} 
\textit{Texto en cursiva.} 
\underline{Texto subrayado.}
\\ \\
Si quieren hacer un salto de línea, es necesario poner \\
Así bajo al siguiente párrafo.

Si quieren dejar una línea en blanco, son 2 saltos de línea \\ \\
Y así quedaría.

\subsubsection{Imágenes}

Para insertar una imagen, simplemente deben hacer:
\\ \\

\begin{figure}[H] %Es importante poner H mayúscula aquí, de este modo la imagen queda fijada en esta posición
    \centering
    \includegraphics[width=0.8\linewidth]{LOGO_UNIVERSIDAD.jpg}
    \caption{Descripción de la imagen}
    \label{fig:LOGO}
\end{figure}

Como se muestra en la Figura \ref{fig:LOGO}, la imagen representa el logo de la universidad.
\\ \\
También puedo poner 2 imágenes juntas:

\begin{figure}[H] %Enfatizo la H mayúscula para que las imágenes queden fijas en esta posición
    \centering
    \begin{minipage}[b]{0.45\linewidth}
        \includegraphics[width=\linewidth]{LOGO_UNIVERSIDAD.jpg}
        \caption{Logo 1}
        \label{fig:logo1}
    \end{minipage}
    \hfill
    \begin{minipage}[b]{0.45\linewidth}
        \includegraphics[width=\linewidth]{LOGO_UNIVERSIDAD.jpg}
        \caption{Logo 2}
        \label{fig:logo2}
    \end{minipage}
\end{figure}

\subsubsection{Fórmulas}

Si quiero insertar una fórmula, el formato más cómodo, a mi parecer, es el siguiente:

\begin{equation}
    \label{eq:Primera-eq}
    \frac{d}{dx} \left( \int_{0}^{x} f(u) \, du \right) = f(x)
\end{equation}

Donde incluso puedo citar la ecuación \ref{eq:Primera-eq} en el texto.
\\ \\
Este formato agrega instantáneamente el número de ecuación.
\\ \\
Como pueden ver, escribir ecuaciones es mucho más cómodo que en Word.

\begin{equation} 
    %No es necesario poner label
    \frac{a}{b} \cdot c = \frac{ac}{b} 
\end{equation}

\subsubsection{Tablas}

Agregar tablas no es complejo; aun así, hay que familiarizarse con el formato.

\begin{table}[H]
    \centering
    \caption{Ejemplo de tabla}
    \label{tab:ejemplo}
    \begin{tabular}{|c|c|c|}
        \hline
        \textbf{Columna 1} & \textbf{Columna 2} & \textbf{Columna 3} \\
        \hline
        Dato 1 & Dato 2 & Dato 3 \\
        \hline
        Dato 4 & Dato 5 & Dato 6 \\
        \hline
    \end{tabular}
\end{table}

Como pueden ver, también hay un label, por lo tanto, las tablas también se pueden citar.

\subsection{Otros tips}

Si hacen sus cálculos en Python, pueden poner su código en LaTeX:

\begin{lstlisting}
    print("Hello World")
\end{lstlisting}

Cualquier cosa, hay documentación de LaTeX en todas partes, por lo tanto, pueden buscar templates o cómo solucionar los distintos problemas.

Recuerden eliminar esta sección para su entrega final.

\end{document}
